An XLSX workbook can contain a small table alongside many empty cell records
and redundant row and column records. Decompressing and parsing these records
adds work. Removing every empty-looking record is
unsafe. An empty cell can preserve a border, override an inherited fill, or
influence automatic row height through its font. A whitespace value can prevent
text from extending into its neighbour. Cleanup must account for how stored records affect the workbook.

The Destatis migration workbook \texttt{wanderungen-2010120217005.xlsx}
illustrates the problem. In \texttt{Tabelle 3.4}, 21 rows contain zeros from
column L to Excel's last column, XFD, totaling 343,833 values. Removing this
hidden content requires explicit removal rules and dependency checks.
Section~\ref{sec:wanderungen} explains the rule and which values it retains.

We designed SheetShrink for use without Windows or Excel, including Python
and R workflows. Direct package processing also avoids Excel's opening cost,
which was high in our experiments.

Cleanup decisions use stored values, styles and references. SheetShrink retains
compact records needed to preserve required workbook behavior. Excel rendering
and calculated row heights are used later to validate the output.

Specialized spreadsheet parsing has demonstrated the importance of decompression,
parsing, memory use and parallel execution for data loading~\cite{sheetreader}.
Our problem adds a transformation obligation. The output remains an editable XLSX
workbook and must retain required data, dependencies and visible presentation
under a stated contract. File size alone is an inadequate measure of success.

Figure~\ref{fig:cleanup-overview} summarizes guarded reductions across cells,
rows, columns and workbook metadata. SheetShrink reuses existing style indices
and compact witnesses to preserve required effects. Native saving can additionally
compact style and shared-string tables.

The performance objective includes the complete user operation. For a sufficiently
bloated workbook, we seek to make cleanup followed by opening the cleaned file
faster than directly opening its original. A single-thread comparison
separates improvements in work performed from gains due to extra processing
resources. Parallel corpus experiments then examine throughput and memory.
